\documentclass[../../main.tex]{subfiles}
\graphicspath{{\subfix{../../image/}}}

\begin{document}

\section{习题}\label{section:ch03-exercises}

在以下习题中，除非另有说明，所有给定函数均假定光滑。

\begin{exercise}\label{exercise:ch03-1}
证明
\begin{align*}
u(x,t,a,b)=a\cdot x-tH(a)+b \qquad(a\in\mathbb{R}^n,\ b\in\mathbb{R})
\end{align*}
是 Hamilton--Jacobi 方程 $u_t+H(Du)=0$ 的一个完全积分。
\end{exercise}

\begin{exercise}\label{exercise:ch03-2}
计算直线族 $x_1+a^2x_2-2a=0$（$a\in\mathbb{R}$）在 $\mathbb{R}^2$ 中的包络，以及平面族 $2a_1x_1+2a_2x_2-x_3+a_1^2+a_2^2=0$（$a_1,a_2\in\mathbb{R}$）在 $\mathbb{R}^3$ 中的包络。画图说明包络的几何意义。
\end{exercise}

\begin{exercise}\label{exercise:ch03-3}
设公式 $G(x,z,a)=0$ 隐式定义函数 $z=u(x,a)$，其中 $x,a\in\mathbb{R}^n$。进一步假定我们可从恒等式
\begin{align*}
\begin{cases}
G(x,u,a)=0,\\
G_{x_i}(x,u,a)+G_z(x,u,a)u_{x_i}=0
\end{cases}\quad(i=1,\cdots,n)
\end{align*}
中消去变量 $a$，解出 $u=u(x)$。
\begin{enumerate}[(a)]
\item 若 $G=\sum_{i=1}^n a_i x_i^2+z^3$，求 $u$ 所满足的 PDE。
\item 在 $\mathbb{R}^{n+1}$ 中，特征刻画所有半径为 1、中心在 $\mathbb{R}^n\times\{z=0\}$ 上的球面的 PDE 是什么？
\end{enumerate}
\end{exercise}

\begin{exercise}\label{exercise:ch03-4}
\begin{enumerate}[(a)]
\item 写出 PDE
\begin{align*}
u_t+b\cdot Du=f \qquad\text{在 }\mathbb{R}^n\times(0,\infty)\text{ 中},
\end{align*}
的特征方程，其中 $b\in\mathbb{R}^n$，$f=f(x,t)$。
\item 用特征 ODE 解 $(*)$，满足初始条件在 $\mathbb{R}^n\times\{t=0\}$ 上 $u=g$。确认你的答案与\eqref{align:ch02-transport-nonhom-sol}一致。
\end{enumerate}
\end{exercise}

\begin{exercise}\label{exercise:ch03-5}
用特征线求解：
\begin{enumerate}[(a)]
\item $x_1u_{x_1}+x_2u_{x_2}=2u$，$u(x_1,1)=g(x_1)$。
\item $x_1u_{x_1}+2x_2u_{x_2}+u_{x_3}=3u$，$u(x_1,x_2,0)=g(x_1,x_2)$。
\item $u u_{x_1}+u_{x_2}=1$，$u(x_1,x_1)=\dfrac{1}{2}x_1$。
\end{enumerate}
\end{exercise}

\begin{exercise}\label{exercise:ch03-6}
给定 $\mathbb{R}^n$ 上光滑向量场 $\mathbf{b}$，令 $\mathbf{x}(s)=\mathbf{x}(s,x,t)$ 解 ODE
\begin{align*}
\begin{cases}
\dot{\mathbf{x}}=\mathbf{b}(\mathbf{x}) & (s\in\mathbb{R}),\\
\mathbf{x}(t)=x.
\end{cases}
\end{align*}
\begin{enumerate}[(a)]
\item 定义 Jacobi 行列式 $J(s,x,t)\triangleq\det D_x\mathbf{x}(s,x,t)$，并推导 Euler 公式：$J_s=(\operatorname{div}\mathbf{b}(\mathbf{x}))J$。
\item 证明 $u(x,t)\triangleq g(\mathbf{x}(0,x,t))J(0,x,t)$ 解
\begin{align*}
\begin{cases}
u_t+\operatorname{div}(u\mathbf{b})=0 & \text{在 }\mathbb{R}^n\times(0,\infty)\text{ 中},\\
u=g & \text{在 }\mathbb{R}^n\times\{t=0\}\text{ 上}.
\end{cases}
\end{align*}
（提示：证明 $\dfrac{\partial}{\partial s}(u(\mathbf{x},s)J)=0$。）
\end{enumerate}
\end{exercise}

\begin{exercise}\label{exercise:ch03-7}
验证 \S3.2.3 断言 (36)：当 $\Gamma$ 在 $x_0$ 附近不平坦时，非特征条件为 $D_pF(p_0,z_0,x_0)\cdot\boldsymbol{\nu}(x_0)\neq 0$。
\end{exercise}

\begin{exercise}\label{exercise:ch03-8}
确认公式 $u=g(x-tF'(u))$（来自 \S3.2.5）为守恒律 $u_t+\operatorname{div}\mathbf{F}(u)=0$ 提供了一个隐式解。
\end{exercise}

\begin{exercise}\label{exercise:ch03-9}
考虑在新容许类
\begin{align*}
\mathcal{A}\triangleq\{\mathbf{w}(\cdot)\in C^2([0,t];\mathbb{R}^n)\mid\mathbf{w}(t)=x\}
\end{align*}
上极小化作用量 $\int_0^t L(\dot{\mathbf{w}}(s),\mathbf{w}(s))\,\mathrm{d}s$ 的问题，其中我们不要求 $\mathbf{w}(0)=y$。
\begin{enumerate}[(a)]
\item 证明极小元 $\mathbf{x}(\cdot)\in\mathcal{A}$ 解 Euler--Lagrange 方程 $-\dfrac{\mathrm{d}}{\mathrm{d}s}(D_vL(\dot{\mathbf{x}}(s),\mathbf{x}(s)))+D_xL(\dot{\mathbf{x}}(s),\mathbf{x}(s))=0$（$0\leqslant s\leqslant t$）。
\item 证明 $D_vL(\dot{\mathbf{x}}(0),\mathbf{x}(0))=0$。
\item 现设 $\mathbf{x}(\cdot)\in\mathcal{A}$ 极小化修正作用量 $\int_0^t L(\dot{\mathbf{w}}(s),\mathbf{w}(s))\,\mathrm{d}s+g(\mathbf{w}(0))$。证明 $\mathbf{x}(\cdot)$ 解通常的 Euler--Lagrange 方程，并确定 $s=0$ 处的边界条件。
\end{enumerate}
\end{exercise}

\begin{exercise}\label{exercise:ch03-10}
若 $H:\mathbb{R}^n\to\mathbb{R}$ 凸，写 $L=H^*$。
\begin{enumerate}[(a)]
\item 设 $H(p)=\dfrac{1}{r}|p|^r$（$1<r<\infty$）。证明 $L(v)=\dfrac{1}{s}|v|^s$，其中 $\dfrac{1}{r}+\dfrac{1}{s}=1$。
\item 设 $H(p)=\dfrac{1}{2}\sum_{i,j=1}^n a_{ij}p_ip_j+\sum_{i=1}^n b_i p_i$，其中 $A=((a_{ij}))$ 对称正定，$b\in\mathbb{R}^n$。计算 $L(v)$。
\end{enumerate}
\end{exercise}

\begin{exercise}\label{exercise:ch03-11}
设 $H:\mathbb{R}^n\to\mathbb{R}$ 凸。我们称 $v$ 属于 $H$ 在 $p$ 处的次微分，记 $v\in\partial H(p)$，若对一切 $r\in\mathbb{R}^n$ 有 $H(r)\geqslant H(p)+v\cdot(r-p)$。证明 $v\in\partial H(p)$ 当且仅当 $p\in\partial L(v)$，当且仅当 $p\cdot v=H(p)+L(v)$，其中 $L=H^*$。
\end{exercise}

\begin{exercise}\label{exercise:ch03-12}
设 $L_1,L_2:\mathbb{R}^n\to\mathbb{R}$ 凸、光滑且超线性。证明
\begin{align*}
\min_{v\in\mathbb{R}^n}(L_1(v)+L_2(v))=\max_{p\in\mathbb{R}^n}(-H_1(p)-H_2(-p)),
\end{align*}
其中 $H_1=L_1^*$，$H_2=L_2^*$。
\end{exercise}

\begin{exercise}\label{exercise:ch03-13}
证明 Hopf--Lax 公式可写为
\begin{align*}
u(x,t)=\min_{y\in\mathbb{R}^n}\left\{tL\!\left(\frac{x-y}{t}\right)+g(y)\right\}
=\min_{y\in B(x,Rt)}\left\{tL\!\left(\frac{x-y}{t}\right)+g(y)\right\},
\end{align*}
其中 $R=\sup_{\mathbb{R}^n}|DH(Dg)|$，$H=L^*$。（这证明了：对具凸 Hamiltonian 且初始函数 $g$ Lipschitz 连续的 Hamilton--Jacobi PDE，传播速度有限。）
\end{exercise}

\begin{exercise}\label{exercise:ch03-14}
设 $E$ 是 $\mathbb{R}^n$ 的闭子集。证明：若 Hopf--Lax 公式可应用于初值问题
\begin{align*}
\begin{cases}
u_t+|Du|^2=0 & \text{在 }\mathbb{R}^n\times(0,\infty)\text{ 中},\\
u=
\begin{cases}
0 & x\in E,\\
+\infty & x\notin E
\end{cases}
& \text{在 }\mathbb{R}^n\times\{t=0\}\text{ 上},
\end{cases}
\end{align*}
则它给出解 $u(x,t)=\dfrac{1}{4t}\operatorname{dist}(x,E)^2$。
\end{exercise}

\begin{exercise}\label{exercise:ch03-15}
给出 \S3.3.3 引理\cref{lemma:ch03-semiconcave2} 证明的全部细节。
\end{exercise}

\begin{exercise}\label{exercise:ch03-16}
设 $u_1,u_2$ 是由 Hopf--Lax 公式给出的两个初值问题
\begin{align*}
\begin{cases}
u_{i,t}+H(Du_i)=0 & \text{在 }\mathbb{R}^n\times(0,\infty)\text{ 中},\\
u_i=g_i & \text{在 }\mathbb{R}^n\times\{t=0\}\text{ 上}
\end{cases}\quad(i=1,2)
\end{align*}
的解。证明 $L^\infty$ 压缩不等式 $\sup_{\mathbb{R}^n}|u_1(\cdot,t)-u_2(\cdot,t)|\leqslant\sup_{\mathbb{R}^n}|g_1-g_2|$（$t>0$）。
\end{exercise}

\begin{exercise}\label{exercise:ch03-17}
证明
\begin{align*}
u(x,t)\triangleq
\begin{cases}
-\dfrac{2}{3}\left(t+\sqrt{3x+t^2}\right) & \text{若 }4x+t^2>0,\\
0 & \text{若 }4x+t^2<0
\end{cases}
\end{align*}
是 $u_t+\left(\dfrac{u^2}{2}\right)_x=0$ 的一个（无界）熵解。
\end{exercise}

\begin{exercise}\label{exercise:ch03-18}
设对一切 $z>0$ 有 $u(x+z)-u(x)\leqslant Ez$。令 $u_\varepsilon=\eta_\varepsilon*u$，证明 $u_{\varepsilon,x}\leqslant E$。
\end{exercise}

\begin{exercise}\label{exercise:ch03-19}
设 $F(0)=0$，$u$ 是守恒律
\begin{align*}
\begin{cases}
u_t+F(u)_x=0 & \text{在 }\mathbb{R}\times(0,\infty)\text{ 中},\\
u=g & \text{在 }\mathbb{R}\times\{t=0\}\text{ 上}
\end{cases}
\end{align*}
的连续积分解，且对每个时刻 $T>0$，$u$ 在 $\mathbb{R}\times[0,T]$ 中具紧支集。证明对一切 $t>0$，$\int_{-\infty}^{\infty}u(\cdot,t)\,\mathrm{d}x=\int_{-\infty}^{\infty}g\,\mathrm{d}x$。
\end{exercise}

\begin{exercise}\label{exercise:ch03-20}
对
\begin{align*}
g(x)=
\begin{cases}
1 & \text{若 }x<-1,\\
0 & \text{若 }-1<x<0,\\
2 & \text{若 }0<x<1,\\
0 & \text{若 }x>1,
\end{cases}
\end{align*}
显式计算
\begin{align*}
\begin{cases}
u_t+\left(\dfrac{u^2}{2}\right)_x=0 & \text{在 }\mathbb{R}\times(0,\infty)\text{ 中},\\
u=g & \text{在 }\mathbb{R}\times\{t=0\}\text{ 上}
\end{cases}
\end{align*}
的唯一熵解。画图记录你的答案，务必说明对一切时刻 $t>0$ 发生了什么。
\end{exercise}

\end{document}